%%
%% Proposal: Dual Numbers for Automatic Differentiation in MScheme
%%
\documentclass[12pt]{report}
\usepackage[margin=1in]{geometry}
\usepackage{booktabs}
\usepackage{listings}
\usepackage{xcolor}
\usepackage{amsmath,amssymb}
\usepackage{enumitem}
\usepackage{hyperref}

\lstdefinelanguage{Scheme}{
  morekeywords={define,lambda,if,cond,else,let,let*,letrec,begin,
    set!,quote,and,or,not,make-rectangular,make-dual,
    dual-real,dual-epsilon,dual?,real-part,imag-part,
    exact->inexact,inexact->exact},
  sensitive=true,
  morecomment=[l]{;},
  morestring=[b]",
}

\lstdefinelanguage{Modula3}{
  morekeywords={MODULE,INTERFACE,PROCEDURE,BEGIN,END,RETURN,IF,THEN,
    ELSE,ELSIF,VAR,CONST,TYPE,IMPORT,FROM,REVEALS,BRANDED,OBJECT,
    METHODS,OVERRIDES,TYPECASE,ISTYPE,NARROW,NEW,RAISES,RAISE,
    TRY,EXCEPT,WITH,FOR,TO,BY,DO,WHILE,LOOP,EXIT,REPEAT,UNTIL,
    CASE,OF,AND,OR,NOT,NIL,TRUE,FALSE,ARRAY,RECORD,REF,REFANY,
    SchemeObject},
  sensitive=true,
  morecomment=[s]{(*}{*)},
  morestring=[b]",
}

\lstset{
  basicstyle=\ttfamily\small,
  keywordstyle=\bfseries,
  commentstyle=\itshape\color{gray},
  stringstyle=\color{brown},
  showstringspaces=false,
  columns=flexible,
  breaklines=true,
}

\title{\bfseries\sffamily Proposal: Dual Numbers for\\Automatic Differentiation in MScheme}
\author{Mika Nystr\"om and Claude (Anthropic)}
\date{March 2026}

\begin{document}
\maketitle
\tableofcontents

%======================================================================
\chapter{Introduction}
%======================================================================

Dual numbers, introduced by William Kingdon Clifford in
1873~\cite{clifford1873}, extend the real numbers with an
infinitesimal element $\varepsilon$ satisfying $\varepsilon^2 = 0$
but $\varepsilon \neq 0$.  A dual number has the form
$a + b\varepsilon$ where $a$ and~$b$ are ordinary numbers.  The
key property is that evaluating any differentiable function at
$a + \varepsilon$ yields $f(a) + f'(a)\varepsilon$---the derivative
appears automatically, with no symbolic differentiation or
finite-difference approximation.

This proposal adds {\tt SchemeDual.T} to MScheme's numeric tower,
following the same architecture as {\tt SchemeComplex.T}.  The
implementation reuses {\tt SchemeNumber} dispatch throughout, so dual
numbers interact correctly with the full tower: exact integers,
rationals, inexact reals, MPFR floats, and complex numbers.

%======================================================================
\chapter{Algebra}
%======================================================================

\section{Arithmetic}

Let $x = a + b\varepsilon$ and $y = c + d\varepsilon$.

\begin{center}
\begin{tabular}{ll}
\toprule
Operation & Result \\
\midrule
$x + y$   & $(a + c) + (b + d)\varepsilon$ \\
$x - y$   & $(a - c) + (b - d)\varepsilon$ \\
$x \cdot y$ & $ac + (ad + bc)\varepsilon$ \\
$x / y$   & $a/c + (bc - ad)/c^2 \cdot \varepsilon$ \\
$-x$      & $-a + (-b)\varepsilon$ \\
\bottomrule
\end{tabular}
\end{center}

\noindent Note that the $bd\varepsilon^2$ term in multiplication
vanishes because $\varepsilon^2 = 0$.

\section{Transcendental Functions}

For any differentiable function $f$:
\[
  f(a + b\varepsilon) = f(a) + b \cdot f'(a) \cdot \varepsilon
\]

\noindent Specific cases:

\begin{center}
\begin{tabular}{ll}
\toprule
Function & Dual result \\
\midrule
$\exp(a + b\varepsilon)$ & $e^a + b e^a \varepsilon$ \\
$\log(a + b\varepsilon)$ & $\log a + b/a \cdot \varepsilon$ \\
$\sin(a + b\varepsilon)$ & $\sin a + b \cos a \cdot \varepsilon$ \\
$\cos(a + b\varepsilon)$ & $\cos a - b \sin a \cdot \varepsilon$ \\
$\tan(a + b\varepsilon)$ & $\tan a + b / \cos^2 a \cdot \varepsilon$ \\
$\sqrt{a + b\varepsilon}$ & $\sqrt{a} + b/(2\sqrt{a}) \cdot \varepsilon$ \\
$(a + b\varepsilon)^n$   & $a^n + n b a^{n-1} \varepsilon$ \\
\bottomrule
\end{tabular}
\end{center}

\section{Interaction with Complex Numbers}

A dual number's parts are {\tt SchemeObject.T}---they can be any
numeric type, including complex.  A dual with complex parts
$(a + bi) + (c + di)\varepsilon$ represents a complex-valued function
and its complex derivative.  The arithmetic rules are unchanged;
the derivative of a complex function is also complex.

There is no ambiguity: ``dual of complex'' and ``complex of dual''
yield isomorphic rings (both have four real components), but we
choose dual-outside because the AD use case is ``evaluate $f$ at
$x + \varepsilon$'' where $x$ may be complex.

\section{Interaction with Exact Arithmetic}

Dual numbers preserve exactness.  If both parts are exact
(integers or rationals), the result is exact:
\begin{lstlisting}[language=Scheme]
(+ (make-dual 1/3 1) (make-dual 1/6 2))  ;=> (make-dual 1/2 3)
\end{lstlisting}
Inexact contagion applies per-part: if either operand's real part
is inexact, the result's real part is inexact (and similarly for
the epsilon part).

%======================================================================
\chapter{Scheme Interface}
%======================================================================

\section{Construction and Accessors}

\begin{description}[style=nextline]
\item[\tt (make-dual \textit{real} \textit{epsilon})]
  Construct a dual number.  If \textit{epsilon} is exactly zero,
  return \textit{real} (demotion, same as complex with zero
  imaginary part).

\item[\tt (dual-real \textit{x})]
  Extract the real part.  For non-dual numbers, returns \textit{x}.

\item[\tt (dual-epsilon \textit{x})]
  Extract the epsilon part.  For non-dual numbers, returns~0.

\item[\tt (dual? \textit{x})]
  Returns {\tt \#t} if \textit{x} is a dual number with nonzero
  epsilon part.
\end{description}

\section{Automatic Differentiation}

\begin{lstlisting}[language=Scheme]
;; Derivative of f at x:
(define (derivative f x)
  (dual-epsilon (f (make-dual x 1))))

;; Example: d/dx (x^2 + sin(x)) at x=3
(derivative (lambda (x) (+ (* x x) (sin x))) 3)
;=> 6 + cos(3) = 5.0099...

;; Works with complex-valued functions:
(derivative (lambda (z) (* z z)) (make-rectangular 1 2))
;=> 2+4i (= 2z at z=1+2i)

;; Works with exact arithmetic:
(derivative (lambda (x) (* x x x)) 1/3)
;=> 1/3 (= 3x^2 at x=1/3)
\end{lstlisting}

\section{Predicates}

\begin{center}
\begin{tabular}{lll}
\toprule
Predicate & Dual number & Rationale \\
\midrule
{\tt number?}  & {\tt \#t} & Duals are numbers \\
{\tt dual?}    & {\tt \#t} & New predicate \\
{\tt complex?} & {\tt \#f} & Not in the complex plane \\
{\tt real?}    & {\tt \#f} & Has nonzero epsilon part \\
{\tt exact?}   & per-part  & Exact iff both parts exact \\
{\tt zero?}    & both zero & $a = 0$ and $b = 0$ \\
\bottomrule
\end{tabular}
\end{center}

\noindent Ordering ({\tt <}, {\tt >}, etc.) on dual numbers
should raise an error, since dual numbers do not have a natural
total order (same rationale as complex).

\section{Display and Reader Syntax}

Dual numbers use the suffix {\tt eps} (for $\varepsilon$), with
{\tt +} or {\tt -} separating the real and epsilon parts.  The
suffix is unambiguous: floating-point exponents use {\tt e} followed
by digits ({\tt 3.14e2}), while the dual suffix is {\tt eps}
followed by end-of-token.

\paragraph{Simple cases (real-valued parts).}

\begin{center}
\begin{tabular}{ll}
\toprule
Dual number & Syntax \\
\midrule
$3 + 4\varepsilon$      & {\tt 3+4eps} \\
$3 - 4\varepsilon$      & {\tt 3-4eps} \\
$0 + 5\varepsilon$      & {\tt 5eps} \\
$0 + \varepsilon$       & {\tt eps} \\
$0 - \varepsilon$       & {\tt -eps} \\
$\frac{1}{3} + 2\varepsilon$ & {\tt 1/3+2eps} \\
$3.14 + 4 \times 10^{10}\varepsilon$ & {\tt 3.14+4.0e10eps} \\
\bottomrule
\end{tabular}
\end{center}

\paragraph{Complex-valued parts.}
When either part is complex (or negative, to avoid {\tt 3+-4eps}),
it is wrapped in parentheses:

\begin{center}
\begin{tabular}{ll}
\toprule
Dual number & Syntax \\
\midrule
$(3+4i) + 5\varepsilon$         & {\tt (3+4i)+5eps} \\
$3 + (4i)\varepsilon$           & {\tt 3+(4i)eps} \\
$(3+4i) + (1+2i)\varepsilon$    & {\tt (3+4i)+(1+2i)eps} \\
$0 + (4i)\varepsilon$           & {\tt (4i)eps} \\
$3 + (-4)\varepsilon$           & {\tt 3+(-4)eps} \\
\bottomrule
\end{tabular}
\end{center}

\noindent The rule for display: if a part is complex, negative, or
itself a dual (if nesting were allowed), wrap it in parentheses.
Otherwise display it bare.

\paragraph{Reader.}
The reader recognizes the {\tt eps} suffix after a numeric literal.
For the {\tt a+b\,eps} form, the {\tt +} or {\tt -} separates two
sub-expressions, with {\tt eps} terminating the epsilon part.
Parenthesized complex parts are parsed as sub-expressions.
Reader support can be deferred to a later phase; {\tt make-dual}
suffices initially.

%======================================================================
\chapter{Modula-3 Implementation}
%======================================================================

\section{Type Definition}

\begin{lstlisting}[language=Modula3]
(* SchemeDual.i3 *)
INTERFACE SchemeDual;
IMPORT SchemeObject, SchemeInt;
FROM Scheme IMPORT E;

TYPE T = SchemeObject.T BRANDED "SchemeDual" OBJECT
  re, eps : SchemeObject.T;
END;

PROCEDURE New(re, eps: SchemeObject.T): SchemeObject.T;
  (* Demotes to re if eps is exactly zero *)

PROCEDURE RealPart(x: T): SchemeObject.T;
PROCEDURE EpsilonPart(x: T): SchemeObject.T;

(* Arithmetic *)
PROCEDURE Add(a, b: T): SchemeObject.T RAISES {E};
PROCEDURE Sub(a, b: T): SchemeObject.T RAISES {E};
PROCEDURE Mul(a, b: T): SchemeObject.T RAISES {E};
PROCEDURE Div(a, b: T): SchemeObject.T RAISES {E};
PROCEDURE Neg(a: T): SchemeObject.T RAISES {E};

PROCEDURE Equal(a, b: T): BOOLEAN RAISES {E};
PROCEDURE Format(x: T): TEXT;

END SchemeDual.
\end{lstlisting}

The {\tt New} constructor demotes to the real part when the epsilon
part is exactly zero, paralleling {\tt SchemeComplex.New} which
demotes when the imaginary part is zero.

\section{Integration with SchemeNumber}

The existing {\tt SchemeNumber} module dispatches on type via
{\tt TYPECASE} and {\tt ISTYPE}.  Adding dual numbers requires
the same pattern used for complex:

\begin{enumerate}[nosep]
\item In {\tt Add}, {\tt Sub}, {\tt Mul}, {\tt Div}, {\tt Neg},
  {\tt Abs}: check for {\tt SchemeDual.T} before the exact/inexact
  dispatch.  When one operand is dual and the other is not,
  promote the non-dual to dual with epsilon = 0.
\item In {\tt Compare}, {\tt Equal}: duals with equal real and
  epsilon parts are equal; ordering raises an error.
\item In {\tt IsExact}, {\tt IsInexact}: check both parts.
\item In {\tt Is}: return {\tt TRUE} for duals.
\item In {\tt Format}: delegate to {\tt SchemeDual.Format}.
\item In {\tt CheckReal}: raise for duals (same as complex).
\end{enumerate}

\section{Integration with SchemePrimitive}

The transcendental functions in {\tt SchemePrimitive.m3} use the
same {\tt TYPECASE} pattern as complex.  For each function $f$,
add a {\tt SchemeDual.T} case that computes
$f(a) + b \cdot f'(a) \cdot \varepsilon$ using the existing
real-valued $f$ and $f'$:

\begin{lstlisting}[language=Modula3]
(* Example: exp on dual *)
| SchemeDual.T(d) =>
  WITH fa  = DoExp(d.re),    (* exp(a) -- recursive call *)
       dfa = fa               (* exp'(a) = exp(a) *)
  DO
    RETURN SchemeDual.New(fa, SchemeNumber.Mul(d.eps, dfa))
  END
\end{lstlisting}

\noindent Since the parts are {\tt SchemeObject.T}, the recursive
call handles the case where {\tt d.re} is complex---{\tt DoExp}
dispatches to {\tt ComplexExp} automatically.

The following is a complete list of every math builtin in
{\tt SchemePrimitive.m3} that must be updated for dual numbers,
grouped by behavior.

\paragraph{Dual-aware (compute $f(a) + b f'(a) \varepsilon$).}

\begin{center}
\small
\begin{tabular}{lll}
\toprule
Builtin & $f(a)$ & $f'(a)$ \\
\midrule
{\tt exp}   & $e^a$          & $e^a$ \\
{\tt log}   & $\log a$       & $1/a$ \\
{\tt sin}   & $\sin a$       & $\cos a$ \\
{\tt cos}   & $\cos a$       & $-\sin a$ \\
{\tt tan}   & $\tan a$       & $1/\cos^2 a$ \\
{\tt asin}  & $\arcsin a$    & $1/\sqrt{1-a^2}$ \\
{\tt acos}  & $\arccos a$    & $-1/\sqrt{1-a^2}$ \\
{\tt atan}  & $\arctan a$    & $1/(1+a^2)$ \\
{\tt sqrt}  & $\sqrt{a}$     & $1/(2\sqrt{a})$ \\
{\tt expt}  & $a^n$          & $n a^{n-1}$ (general: $a^n \log a$) \\
{\tt abs}   & $|a|$          & $\mathrm{sign}(a)$ (real duals only) \\
{\tt tanh}  & $\tanh a$      & $1/\cosh^2 a$ \\
\bottomrule
\end{tabular}
\end{center}

\paragraph{Dual arithmetic (already covered by SchemeNumber).}
These are handled in {\tt SchemeNumber.m3}, not per-builtin:
{\tt +}, {\tt -}, {\tt *}, {\tt /}, unary {\tt -}.

\paragraph{Predicates (check/dispatch on dual type).}

\begin{center}
\begin{tabular}{ll}
\toprule
Builtin & Behavior for dual \\
\midrule
{\tt number?}  & {\tt \#t} \\
{\tt dual?}    & {\tt \#t} (new predicate) \\
{\tt complex?} & {\tt \#f} \\
{\tt real?}    & {\tt \#f} \\
{\tt rational?} & {\tt \#f} \\
{\tt integer?} & {\tt \#f} \\
{\tt exact?}   & exact iff both parts exact \\
{\tt inexact?} & inexact iff either part inexact \\
{\tt zero?}    & both parts zero \\
{\tt positive?} & error (not ordered) \\
{\tt negative?} & error (not ordered) \\
{\tt even?}    & error \\
{\tt odd?}     & error \\
\bottomrule
\end{tabular}
\end{center}

\paragraph{Real-only (error on dual, via {\tt CheckReal}).}
These operations are not defined for dual numbers and should
raise ``not a real number'':

\begin{center}
\begin{tabular}{l}
\toprule
Builtins \\
\midrule
{\tt <\quad >\quad <=\quad >=\quad min\quad max} \\
{\tt floor\quad ceiling\quad truncate\quad round} \\
{\tt quotient\quad remainder\quad modulo\quad gcd\quad lcm} \\
\bottomrule
\end{tabular}
\end{center}

\paragraph{Equality.}
{\tt =} returns {\tt \#t} when both real and epsilon parts are equal.
{\tt eqv?} and {\tt equal?} use identity/structural comparison as usual.

\paragraph{Conversion.}

\begin{center}
\begin{tabular}{ll}
\toprule
Builtin & Behavior \\
\midrule
{\tt exact->inexact} & convert both parts \\
{\tt inexact->exact} & convert both parts \\
{\tt number->string} & delegate to {\tt SchemeDual.Format} \\
\bottomrule
\end{tabular}
\end{center}

\paragraph{Accessors (new).}

\begin{center}
\begin{tabular}{ll}
\toprule
Builtin & Behavior \\
\midrule
{\tt dual-real}    & real part; identity for non-duals \\
{\tt dual-epsilon} & epsilon part; 0 for non-duals \\
{\tt make-dual}    & constructor (demotes if epsilon = 0) \\
\bottomrule
\end{tabular}
\end{center}

\section{Interop (SchemeModula3Types)}

{\tt SchemeDual.T} is a traced M3 reference type, so it passes
through the Scheme--M3 boundary like any other {\tt REFANY}.
Modula-3 interfaces that declare parameters of type
{\tt SchemeDual.T} receive the dual object directly, and can
access its {\tt re} and {\tt eps} fields from M3 code.

{\tt SchemeModula3Types} should add {\tt SchemeDual\_T} as a
recognized interop type (parallel to {\tt Mpz.T} and {\tt Mpfr.T}),
with {\tt ToScheme}/{\tt ToModula} converters and
an entry in the {\tt Type} enumeration.  The {\tt interop-roundtrip}
primitive should support {\tt 'SchemeDual\_T}.

For M3 formal types that expect real numbers ({\tt INTEGER},
{\tt LONGREAL}, etc.), dual arguments should raise an error,
same as complex.

\section{Modula-3 API for Complex and Dual}

{\tt Mpz.T} is a pre-existing M3 library with its own rich API
(200+ operations).  {\tt SchemeComplex.T} and {\tt SchemeDual.T}
are different: they originate in the Scheme numeric tower, so
their M3 interfaces must be designed to be useful from both sides.

The complex transcendental functions ({\tt ComplexExp},
{\tt ComplexLog}, {\tt ComplexSin}, etc.) are currently static
procedures inside {\tt SchemePrimitive.m3}.  These should move
into {\tt SchemeComplex.i3/m3} as public exports, and the
corresponding dual versions into {\tt SchemeDual.i3/m3}:

\begin{lstlisting}[language=Modula3]
(* SchemeComplex.i3 — public M3 API *)
PROCEDURE Exp(c: T): SchemeObject.T RAISES {E};
PROCEDURE Log(c: T): SchemeObject.T RAISES {E};
PROCEDURE Sin(c: T): SchemeObject.T RAISES {E};
PROCEDURE Cos(c: T): SchemeObject.T RAISES {E};
PROCEDURE Sqrt(c: T): SchemeObject.T RAISES {E};
(* ... and similarly for SchemeDual.i3 *)
\end{lstlisting}

\noindent With these as public interfaces, M3 code can call
complex and dual math directly:

\begin{lstlisting}[language=Modula3]
VAR z := SchemeComplex.New(SchemeInt.FromI(3),
                           SchemeInt.FromI(4));
VAR w := SchemeComplex.Exp(z);   (* e^(3+4i) *)
\end{lstlisting}

\noindent {\tt SchemePrimitive.m3} then dispatches the Scheme
builtins ({\tt exp}, {\tt sin}, etc.) to these module procedures
rather than inlining the formulas.

{\tt sstubgen} can generate Scheme stubs for these interfaces
if requested (e.g., via {\tt SchemeStubs("SchemeComplex")} in
a package's m3makefile), making the operations callable from
Scheme as {\tt (SchemeComplex.Exp z)}.  This falls out naturally
from the existing stub generation mechanism---no special support
is needed.

The full public API for each module:

\begin{center}
\small
\begin{tabular}{ll}
\toprule
Category & Procedures \\
\midrule
Construction   & {\tt New}, {\tt MakeRectangular}/{\tt MakePolar} (complex) \\
               & {\tt New} (dual) \\
Accessors      & {\tt RealPart}, {\tt ImagPart}, {\tt Magnitude}, {\tt Angle} (complex) \\
               & {\tt RealPart}, {\tt EpsilonPart} (dual) \\
Arithmetic     & {\tt Add}, {\tt Sub}, {\tt Mul}, {\tt Div}, {\tt Neg} \\
Comparison     & {\tt Equal} \\
Transcendental & {\tt Exp}, {\tt Log}, {\tt Sin}, {\tt Cos}, {\tt Tan} \\
               & {\tt Asin}, {\tt Acos}, {\tt Atan}, {\tt Sqrt} \\
Query          & {\tt IsZero} \\
Display        & {\tt Format} \\
\bottomrule
\end{tabular}
\end{center}

%======================================================================
\chapter{Compiler Support}
%======================================================================

\section{Constant Embedding}

The compiler's {\tt constant-ref} and {\tt constant-to-m3} need
a dual case, parallel to the complex case:

\begin{lstlisting}[language=Scheme]
;; In constant-ref / constant-to-m3:
((dual? val)
 (string-append "SchemeDual.New("
                (constant-ref (dual-real val) ctx) ", "
                (constant-ref (dual-epsilon val) ctx) ")"))
\end{lstlisting}

\section{{\tt form->string}}

The {\tt form->string} serializer needs to handle duals for the
{\tt loadEvalText} fallback path:

\begin{lstlisting}[language=Scheme]
((dual? form)
 (string-append "(make-dual "
                (form->string (dual-real form)) " "
                (form->string (dual-epsilon form)) ")"))
\end{lstlisting}

\section{Inline Primitives}

The {\tt number?} predicate already compiles to
{\tt SchemeNumber.Is()}, which will return {\tt TRUE} for duals
once {\tt SchemeNumber.Is} is updated.  No compiler change needed.

The {\tt dual?} predicate can be inlined as
{\tt ISTYPE(x, SchemeDual.T)} in {\tt gen-inline-expr}.

Arithmetic ({\tt +}, {\tt -}, {\tt *}, {\tt /}) goes through
{\tt SchemePrimitive.NumericAdd} etc., which delegate to
{\tt SchemeNumber.Add}---no compiler change needed.

\section{IMPORT Line}

Add {\tt SchemeDual} to the generated M3 module's IMPORT list
(alongside the existing {\tt SchemeComplex}, {\tt SchemeRational},
etc.).

\section{Unboxing}

Numeric unboxing (currently disabled) assumes {\tt LONGREAL}
representation.  Duals cannot be unboxed to {\tt LONGREAL}.
The existing guard that disables unboxing for the numeric tower
covers this case; no change needed.

%======================================================================
\chapter{Testing}
%======================================================================

A test file {\tt test-dual.scm} should cover:

\begin{enumerate}[nosep]
\item Basic arithmetic: {\tt +}, {\tt -}, {\tt *}, {\tt /}, unary {\tt -}
\item Demotion: {\tt (make-dual 5 0)} $\to$ {\tt 5}
\item Predicates: {\tt dual?}, {\tt number?}, {\tt real?}, {\tt exact?}
\item Transcendentals: {\tt exp}, {\tt log}, {\tt sin}, {\tt cos},
  {\tt tan}, {\tt sqrt}, {\tt expt}, {\tt asin}, {\tt acos}, {\tt atan}
\item Automatic differentiation: verify $f'(x)$ for known functions
\item Exact arithmetic: rational-valued duals
\item Complex-valued duals: derivatives of complex functions
\item Error cases: ordering, {\tt floor}, {\tt positive?}, etc.
\item {\tt exact->inexact}, {\tt inexact->exact} on duals
\item {\tt number->string} on duals
\item Interop: {\tt REFANY} roundtrip
\item Loops with dual accumulators (compiler acceptance)
\end{enumerate}

A compiler version {\tt test-cn-dual.scm} in the numtest suite
verifies identical behavior compiled vs.\ interpreted.

%======================================================================
\chapter{Implementation Plan}
%======================================================================

\begin{center}
\begin{tabular}{clll}
\toprule
Step & Task & Files & Effort \\
\midrule
1 & {\tt SchemeDual.i3/m3} & new & 1 hour \\
2 & {\tt SchemeNumber.m3} dispatch & modify & 1 hour \\
3 & {\tt SchemePrimitive.m3} transcendentals & modify & 2 hours \\
4 & {\tt SchemeInputPort.m3} reader (optional) & modify & 1 hour \\
5 & {\tt compiler.scm} constants/serialization & modify & 30 min \\
6 & {\tt test-dual.scm} + {\tt test-cn-dual.scm} & new & 2 hours \\
\bottomrule
\end{tabular}
\end{center}

\noindent Steps 1--3 follow the same pattern as the complex number
implementation.  Step~4 (reader syntax) can be deferred.
Step~5 is mechanical---same changes as were made for complex.
Step~6 should be written first (acceptance tests) before
implementing, as was done for the complex math builtins.

{
\frenchspacing
\begin{thebibliography}{99}

\bibitem{clifford1873} William Kingdon Clifford.
  Preliminary sketch of biquaternions.
  {\it Proceedings of the London Mathematical Society,}
  s1-4(1):381--395, 1873.

\end{thebibliography}
}

\end{document}
